%%%%%%%%%%  scan idx 111  |  folio 100  %%%%%%%%%%

\medskip
\noindent\textbf{Teorema 4.5.}
Sea $X_0$ un esquema de tipo finito sobre $k$ para el cual se cumple
una de las condiciones siguientes:

\smallskip
\noindent(i) $X_0$ es propio sobre $k$, o bien

\noindent(ii) $X_0$ es afín y
\[
X_0^{\mathrm{sing}}
=\{x\in X_0\mid X_0\text{ no es liso sobre }k\text{ en }x\}
\]
es finito.

\smallskip
\noindent\textbf{(a)}
La categoría $\underline M_{X_0}$ admite un objeto formal versal
sobre $\underline C_\Lambda$, llamado deformación formal versal de
$X_0$. Si $(R,\mathfrak X)$ es una deformación formal versal de
$X_0$, entonces para todo
$(S,\mathfrak Y)\in\operatorname{Ob}(\widehat{\underline M}_{X_0})$
existe una aplicación
\[
(\varphi,\Phi):(R,\mathfrak X)\longrightarrow(S,\mathfrak Y)
\]
en $\widehat{\underline M}_{X_0}$ tal que el diagrama siguiente es
cartesiano:
\[
\begin{tikzcd}[column sep=huge,row sep=large]
\mathfrak Y \arrow[r,"\Phi"] \arrow[d]
& \mathfrak X \arrow[d] \\
\operatorname{Spf}(S)
 \arrow[r,"\operatorname{Spf}(\varphi)"']
& \operatorname{Spf}(R).
\end{tikzcd}
\]
Además, una deformación formal versal $(R,\mathfrak X)$ de $X_0$
prorepresenta el funtor $[\underline M_{X_0}]$ si y solo si, para toda
aplicación sobreyectiva $R'\longrightarrow R$ en
$\underline C_\Lambda$ y toda deformación $(R',X')$ de $X_0$ sobre
$R'$, el homomorfismo
\[
\operatorname{Aut}_{R'}(X'\mid X_0)
\longrightarrow
\operatorname{Aut}_R(X'\otimes_{R'}R\mid X_0)
\]
es sobreyectivo.

\smallskip
\noindent\textbf{(b)}
Si $X_0$ es propio sobre $k$, entonces $\underline M_{X_0}$ es
pseudorrepresentable de manera lisa sobre $\underline C_\Lambda$.

%%%%%%%%%%  scan idx 112  |  folio 101  %%%%%%%%%%

\medskip
\noindent\textit{Demostración.}
\textbf{(a)} Por el Corolario 4.3, $\underline M_{X_0}$ es una
categoría cofibrada homogénea en grupoides sobre
$\underline C_\Lambda$. Por 1.11 y 2.9(b), basta, por tanto, demostrar
que $[\underline M_{X_0}(k[\varepsilon])]$ tiene dimensión finita
sobre $k$.

Consideremos la sucesión exacta
\begin{equation}\tag{*}
0\longrightarrow H^1(X_0,\underline D^0_{X_0})
\longrightarrow[\underline M_{X_0}(k[\varepsilon])]
\longrightarrow H^0(X_0,\underline D^1_{X_0}).
\end{equation}
% Canon: SGA-SRC-001245.
Si $X_0$ es propio sobre $k$, tanto
$H^1(X_0,\underline D^0_{X_0})$ como
$H^0(X_0,\underline D^1_{X_0})$ tienen dimensión finita, pues los
$\underline D^i_{X_0}$ son coherentes para todo $i\geq0$ por
\sisi{2.13(a)}{3.15(a)}.

Supongamos ahora que $X_0$ es afín y que
$X_0^{\mathrm{sing}}$ es finito. El haz
$\underline D^1_{X_0}$ es coherente y
\[
\operatorname{Supp}\underline D^1_{X_0}
\subset X_0^{\mathrm{sing}}
\]
por \sisi{2.13(a)}{3.15(a)}. Por consiguiente,
$H^0(X_0,\underline D^1_{X_0})$ tiene dimensión finita sobre $k$,
mientras que $H^1(X_0,\underline D^0_{X_0})=0$ porque $X_0$ es afín.
En ambos casos, la sucesión exacta $(*)$ muestra que
$[\underline M_{X_0}(k[\varepsilon])]$ tiene dimensión finita sobre
$k$. La segunda afirmación de la parte (a) se deduce de 2.7.

\smallskip
\noindent\textbf{(b)} Si $X_0$ es propio, entonces
\[
\operatorname{Aut}(X_0\otimes_k k[\varepsilon]\mid X_0)
=H^0(X_0,\underline D^0_{X_0})
\]
tiene dimensión finita, y la afirmación se deduce de 2.13.
\hfill$\square$
